\section{Actual old-branch superlevels avoid born boundary and preserve marked cosupport}
\label{sec:hironaka-normal-crossings-boundary-avoidance}

Here an actual boundary-adapted pair $(H,E)$ consists of an intrinsic
normal-crossing carrier $H$ and a genuine finite SNC list $E$ with
simultaneous coordinates: the local coordinates of the distinct $H$
branches are disjoint from those of every boundary member through the
point. The members of $E$ may be the actual born exceptional members;
no initial merely-NC boundary is silently made SNC.

\begin{theorem}[Every intrinsic old-branch superlevel is dense away from the actual boundary]
\label{hir:general-nc-actual-branch-superlevel-off-boundary-density}
\lean{Hironaka.HironakaNormalCrossingsBoundaryAvoidance.branch_superlevel_subset_closure_off_boundary}
For every actual boundary-adapted pair $(H,E)$ and every natural $r$,
the intrinsic branch-count superlevel $D_r(H)$ is contained in the
closure of $D_r(H)\setminus |E|$. In particular the entire actual
double locus is approximated by actual double-locus points outside
the born boundary.
\uses{hir:general-nc-simultaneous-original-boundary,hir:general-nc-exact-nearby-branch-count}
\leanok
\end{theorem}
\begin{proof}
Fix a point of the superlevel and genuine simultaneous centred
coordinates. Let $S$ be its actual old-branch coordinate set. The
intrinsic branch-count formula gives $r\leq |S|$. Take the vector
whose coordinates in $S$ are zero and whose other coordinates are
one. Multiply it by a small nonzero complex parameter and apply the
actual inverse chart. The resulting actual points tend to the fixed
point and, eventually, belong to the genuine chart target and any
prescribed neighbourhood.

Every old-branch coordinate stays zero, so the exact original-chart
branch-count formula gives the same count $|S|$ at these points.
Every boundary member through the fixed point has a singleton
coordinate outside $S$; that coordinate is the nonzero parameter,
so none of these boundary members contains the perturbed point.
The finitely many boundary members not through the fixed point are
closed, so one further genuine neighbourhood avoids all of them.
Thus every neighbourhood meets $D_r(H)\setminus |E|$, proving the
asserted closure inclusion. Empty boundary lists and zero-dimensional
cases are covered by the same finite coordinate argument.
\end{proof}

\begin{theorem}[Actual closed current cosupport propagates across the born boundary at every mark]
\label{hir:general-nc-actual-closed-cosupport-off-boundary-propagation}
\lean{Hironaka.HironakaNormalCrossingsBoundaryAvoidance.branch_superlevel_subset_cosupp_of_off_boundary}
Let an actual current marked datum have coherent ideal and closed
smooth ambient carrier, and let $(H,E)$ be adapted with its actual
boundary list. If its actual cosupport contains $D_r(H)\setminus|E|$,
then it contains the whole $D_r(H)$. The marking is arbitrary and the
same current ideal, carrier and mark are retained.
\uses{hir:general-nc-actual-branch-superlevel-off-boundary-density,hir:order-power-along}
\leanok
\end{theorem}
\begin{proof}
The actual coherent-ideal cosupport on the actual closed smooth
carrier is closed at the given marking. It contains the assumed
actual off-boundary superlevel, hence contains its closure. The
preceding genuine density theorem places the entire superlevel
inside that same closure. Compose these inclusions. No extra ideal
power or exceptional-factor premise is introduced.
\end{proof}

\begin{theorem}[A genuine highest-old-branch blowup preserves double-locus cosupport for every marking]
\label{hir:general-nc-actual-highest-step-any-mark-double-cosupport}
\lean{Hironaka.HironakaNormalCrossingsBoundaryAvoidance.highest_step_doubleLocus_subset_controlled_cosupport}
For a whole ambient well-formed marked datum and an actual adapted
carrier $H$, suppose its actual original double locus is contained in
the original marked cosupport. Every genuine blowup of a highest
old-branch superlevel of count at least two is an actual admissible
step for this same datum. The actual strict old-carrier double locus
is contained in the literal marking-$d$ controlled cosupport of the
same original ideal after this step, for every marking $d$.
\uses{hir:general-nc-actual-closed-cosupport-off-boundary-propagation,hir:general-nc-boundary-transform-original-count-step,hir:general-nc-higher-count-over-image}
\leanok
\end{theorem}
\begin{proof}
The genuine highest old-branch centre is closed and smooth and has
actual normal crossings with the given SNC boundary. Its count is
at least two, so the original double-locus hypothesis puts this very
centre in the actual original cosupport. It is therefore admissible
for the same original marked datum.

Actual highest-stratum blowup geometry supplies the transformed
adapted pair, including the literal new exceptional member. Outside
its whole actual boundary, a point is outside that exceptional
member, hence projects off the actual centre. The exact original-image
branch-count bound sends an actual new double-locus point to an
actual original double-locus point. The hypothesis puts its image
in the original cosupport. The genuine off-centre controlled-transform
law then puts the point in the actual new cosupport. Actual admissible
transform well-formedness makes this current cosupport closed.
Apply the preceding propagation theorem to the same transformed
adapted pair. This gives the entire new double-locus inclusion,
including exceptional points. Finally the literal whole-ambient
transform dictionary identifies the current ideal and mark with the
actual marking-$d$ controlled transform of the original ideal.
No marking-one restriction, stronger centre-power condition, or
full initial-NC resolution is asserted.
\end{proof}
